
\documentclass[11pt, a4paper]{article}
\usepackage[utf8]{inputenc}
\usepackage{amsmath}
\usepackage{amssymb}
\usepackage{graphicx}
\usepackage{geometry}
\usepackage{hyperref}
\usepackage{booktabs}
\usepackage{siunitx}
\usepackage{amsthm}

\geometry{left=2.5cm, right=2.5cm, top=2.5cm, bottom=2.5cm}

\title{\textbf{Mass Generation via Topological Shielding Gates in Swirl String Theory}\\
\large Rational-Concordance Constraints (Double-Twist) and an SEMF Bridge to Atomic Masses}
\author{Auteur Naam \\ \textit{Swirl String Theory Research Group}}
\date{\today}

% ==========================================================
% SST house macros (compile-safe; match Canon naming)
% ==========================================================
% context-aware swirl arrows (robust sizing)
\newcommand{\swirlarrow}{%
  \mathchoice{\mkern-2mu\scriptstyle\boldsymbol{\circlearrowleft}}%
             {\mkern-2mu\scriptstyle\boldsymbol{\circlearrowleft}}%
             {\mkern-2mu\scriptscriptstyle\boldsymbol{\circlearrowleft}}%
             {\mkern-2mu\scriptscriptstyle\boldsymbol{\circlearrowleft}}%
}
\newcommand{\swirlarrowcw}{%
  \mathchoice{\mkern-2mu\scriptstyle\boldsymbol{\circlearrowright}}%
             {\mkern-2mu\scriptstyle\boldsymbol{\circlearrowright}}%
             {\mkern-2mu\scriptscriptstyle\boldsymbol{\circlearrowright}}%
             {\mkern-2mu\scriptscriptstyle\boldsymbol{\circlearrowright}}%
}
\newcommand{\vswirl}{\mathbf{v}_{\swirlarrow}}
\newcommand{\vswirlcw}{\mathbf{v}_{\swirlarrowcw}}
\newcommand{\vnorm}{\lVert \vswirl \rVert}
\newcommand{\rhoF}{\rho_{\!f}}   % effective fluid density
\newcommand{\rhoE}{\rho_{\!E}}   % swirl energy density
\newcommand{\rhoM}{\rho_{\!m}}   % mass-equivalent density (=rhoE/c^2)
\providecommand{\rc}{r_c}

\begin{document}

\maketitle

\begin{abstract}
We present a self-contained effective mass model in which invariant rest mass is identified with a localized swirl-energy integral and modulated by topology through a compact, testable kernel. The central deliverable is a single factorized master equation
\[
M(T)=\Big(\frac{4}{\alpha}\Big)^{S(T)}\,k(T)^{-3/2}\,\phi^{-g_4(T)}\,n(T)^{-1/\phi}\;\Big(M_0\,L_{\mathrm{ideal}}(T)\Big),
\]
where \(L_{\mathrm{ideal}}\) is a dimensionless ropelength proxy, \(k\) is a spectral/complexity index, \(n\) is component count, and \(g_4\) is the \emph{slice genus} (a concordance-invariant genus measure). We explicitly constrain the kernel by requiring robustness under (rational) concordance moves and, as a concrete genus-one testbed, we incorporate the double-twist family \(K_{m,n}\) and its slice criterion, which provides a sharp on/off benchmark for the genus factor \cite{Lee2025DoubleTwist}. Non-original ingredients (ropelength energetics, discrete scale invariance, concordance constraints) are cited; SST-specific identifications are labeled as postulates.
\par\smallskip
\textbf{New in SST-59 (hard-formalized):} (i) we replace the narrative ``shielding'' discussion by a \emph{binary shielding gate}
\(\mathsf{G}(T)\in\{0,1\}\) that decides whether the amplification \((4/\alpha)\) is active; (ii) we incorporate rational concordance of
double-twist knots \(K_{m,n}\) as a \emph{hard constraint} on \(\mathsf{G}\) via a theorem-level characterization of rational sliceness; and
(iii) we add a bridge model to atomic masses by subtracting an SEMF-like binding functional that is explicitly compatible with the master kernel.
\end{abstract}

\paragraph{Scope \& status.}
This paper (Part~II) formalizes \emph{topological admissibility constraints} on topology-indexed mass kernels used in the effective mass functional program. The goal is not to postulate new forces but to restrict the space of allowed kernel assignments by computable 4D-topological criteria.

\paragraph{Assumptions (Euler--Kelvin + topological sectoring).}
\begin{itemize}
    \item \textbf{Same dynamical regime as Part~I}: incompressible, inviscid, barotropic; no reconnection.
    \item \textbf{Topological sector decomposition}: configurations are partitioned into equivalence classes $T$ (knot/link types plus admissible deformations) with well-defined invariants.
    \item \textbf{Kernel ansatz}: a dimensionless kernel $\mathcal{K}(T)$ is permitted only when $T$ satisfies a stated admissibility criterion; otherwise it is excluded from the model class.
    \item \textbf{No continuous Yukawa-like tuning}: only discrete indices/invariants label admissible kernel structure in this part.
\end{itemize}

\paragraph{What is claimed vs.\ not claimed.}
\begin{itemize}
    \item \textbf{Claimed:} A binary admissibility map $\mathsf{G}(T)\in\{0,1\}$ can be defined from concordance-based (4D-topological) constraints and used as a \emph{selection rule} that restricts the allowed kernel classes $\mathcal{K}(T)$.
    \item \textbf{Not claimed:} We do not claim a proven microscopic necessity that a specific constant (e.g.\ $4/\alpha$) is ``dictated'' by a single invariant; rather, we present a falsifiable constraint program: certain kernel normalizations/ans\"atze are allowed or excluded by computable topological criteria.
\end{itemize}

\paragraph{Observable definition (what is measured).}
The observable output of Part~II is a set of \emph{testable admissibility predictions}:
given a configuration class $T$, the model predicts either
\begin{equation}
\mathsf{G}(T)=1 \ \ (\text{admissible}) \quad \text{or} \quad \mathsf{G}(T)=0 \ \ (\text{excluded}),
\label{eq:obs_partII_gate}
\end{equation}
together with the induced restriction on the mass-kernel assignment $\mathcal{K}(T)$.
These predictions become measurable in Part~III by checking whether only admissible classes can consistently reproduce mass hierarchies under a single fixed normalization, and by identifying outliers as violations of the admissibility criterion rather than as fit freedom.

\section{Introduction}
Standard Model particle masses are currently free parameters determined by Yukawa couplings to the Higgs field. Swirl String Theory (SST) proposes an alternative origin: mass is the invariant energy of a topological defect (vortex) in a continuous medium.

In this framework, the rest mass $M$ is defined by the volume integral of the swirl energy density $\rho_E$:
\begin{equation}
    M = \frac{1}{c^2} \iiint_V \left( \frac{1}{2}\rho_{\text{core}} \vnorm^2 + P_{\text{top}} \right)\, dV
\end{equation}
This paper derives the scaling laws that map knot topology ($T$) to physical mass, identifying three key mechanisms: geometric ropelength, spectral suppression, and vacuum coupling amplification.

\section{The SST Master Equation}
We propose the following unified mass formula for a stable topological defect $T$:

\begin{equation}
    M(T) =
    \underbrace{\left( \frac{4}{\alpha} \right)^{\mathsf{G}(T)}}_{\text{Vacuum Coupling (gate)}}
    \cdot \underbrace{n^{-1/\phi}}_{\text{Coherence}}
    \cdot \underbrace{k^{-3/2}}_{\text{Spectral}}
    \cdot \underbrace{\left( M_0 \cdot L_{\text{ideal}}(T) \right)}_{\text{Geometric Base}}
\end{equation}

Where:
\begin{itemize}
    \item $M_0$: The fundamental vortex mass unit per unit dimensionless length.
    \item $L_{\text{ideal}}$: The ropelength of the ideal knot configuration.
    \item $k$: The dominant spectral mode (symmetry order) of the knot.
    \item $n$: The number of linked components.
    \item $\alpha^{-1}$: The fine-structure constant ($\approx 137.036$).
    \item $\mathsf{G}(T)\in\{0,1\}$: the \emph{shielding gate}.  \(\mathsf{G}(T)=0\) means ``perfect shielding'' and \((4/\alpha)\) is \emph{off};
    \(\mathsf{G}(T)=1\) means shielding is broken and \((4/\alpha)\) is \emph{on}.  Section~\ref{sec:gate} makes this formal.
    \item $\phi$: The Golden Ratio ($\approx 1.618$), governing topological coherence.
\end{itemize}

\section{A formal shielding gate from rational sliceness}\label{sec:gate}
\subsection{Why a gate (not a narrative)}
The master kernel uses a large discrete amplification scale \((4/\alpha)\).  To keep the model checkable, the activation of this factor must be
\emph{decidable} from topology/invariants rather than prose.  We therefore define a binary gate \(\mathsf{G}(T)\).

\subsection{Rational concordance background (definitions)}
\textbf{Convention.} We work in the smooth category.

\begin{definition}[Rationally slice knot]\label{def:rationalslice}
A knot \(K\subset S^3\) is \emph{rationally slice} if it bounds a smoothly embedded disk \(D\subset W\) in a smooth \(4\)-manifold \(W\) with
\(\partial W=S^3\) and \(H_*(W;\mathbb{Q})\cong H_*(B^4;\mathbb{Q})\).
\end{definition}

\begin{definition}[Rational concordance group \(\mathcal{C}_{\mathbb{Q}}\)]
Two knots \(K_0,K_1\subset S^3\) are \emph{rationally concordant} if \(K_0\#-K_1\) is rationally slice.
The set of equivalence classes under connected sum forms an abelian group \(\mathcal{C}_{\mathbb{Q}}\) with identity the class of rationally slice knots.
\end{definition}

\subsection{Shielding gate (axiomatically tied to \(\mathcal{C}_{\mathbb{Q}}\))}
\begin{definition}[Shielding gate]\label{def:gate}
Define the shielding gate
\[
\mathsf{G}(T)=
\begin{cases}
0,& \text{if the knot/link type assigned to }T\text{ is rationally slice (i.e. }[T]=0\text{ in }\mathcal{C}_{\mathbb{Q}}\text{)}\\
1,& \text{otherwise.}
\end{cases}
\]
Equivalently, \(\mathsf{G}\) turns \((4/\alpha)\) \emph{off} exactly on the rationally-slice (``perfectly shielded'') topological sector.
\end{definition}

\paragraph{Operational consequence.}
Given \(\mathsf{G}(T)\), the mass kernel is now strictly:
\[
M(T)=\left(\frac{4}{\alpha}\right)^{\mathsf{G}(T)}\,n^{-1/\phi}\,k^{-3/2}\,M_0\,L_{\mathrm{ideal}}(T).
\]

\section{Double-twist knots as a hard constraint on the gate}\label{sec:doubletwist}
\subsection{Definition of the family \(K_{m,n}\)}
\begin{definition}[Double twist knot \(K_{m,n}\)]\label{def:Kmn}
Let \(K_{m,n}\) denote the double-twist knot with integer parameters \(m,n\), defined by the standard two-box double-twist diagram with
\(m\) full twists in one twist region and \(n\) full twists in the other (sign convention as in \cite{Lee2025DoubleTwist}).
\end{definition}

\subsection{Exact rational sliceness characterization}
\begin{theorem}[Rational sliceness of \(K_{m,n}\)]\label{thm:rationalsliceKmn}
The double twist knot \(K_{m,n}\) is rationally slice if and only if
\[
mn=0 \quad\text{or}\quad n=-m \quad\text{or}\quad n=-m\pm 1.
\]
\end{theorem}
\noindent\emph{Source:} Theorem~1.1 of \cite{Lee2025DoubleTwist}.  Earlier work established the ``if'' direction for
subfamilies (e.g. \(K_{m,-m}\) rationally slice); the converse is the hard content used here as a constraint.

\begin{corollary}[Explicit shielding gate on \(K_{m,n}\)]\label{cor:gateKmn}
For the double-twist family,
\[
\mathsf{G}(K_{m,n})=
\begin{cases}
0,& mn=0\ \text{or}\ |m+n|\le 1,\\
1,& \text{otherwise.}
\end{cases}
\]
\end{corollary}

\paragraph{Model use.}
Whenever a particle assignment uses a double-twist sector \(K_{m,n}\), \(\mathsf{G}\) is no longer adjustable: it is fixed by
Corollary~\ref{cor:gateKmn}. This converts ``shielding'' into a falsifiable topological decision rule.

\section{Bridge to atomic masses via an SEMF-compatible binding functional}\label{sec:semf}
\subsection{What SEMF is (and why it is a bridge)}
The semi-empirical mass formula (SEMF, Bethe--Weizs\"acker liquid-drop model) approximates nuclear binding energy using volume/surface/Coulomb/asymmetry/pairing
terms.  In SST-59 it functions as a \emph{phenomenological bridge}: it corrects the sum of constituent SST rest masses by a binding-energy subtraction,
without modifying the fundamental particle kernel.

\subsection{Atomic mass master equation (kernel + binding subtraction)}
Let \(A=Z+N\).  Define the atomic mass (isotope) as
\begin{equation}
M_{\mathrm{atom}}(Z,N)\;=\; Z\,M_p^{\mathrm{SST}} + N\,M_n^{\mathrm{SST}} + Z\,M_e^{\mathrm{SST}}
\;-\;\frac{E_{\mathrm{bind}}^{\mathrm{bridge}}(Z,N)}{c^2},
\label{eq:MatomBridge}
\end{equation}
where each \(M_{p,n,e}^{\mathrm{SST}}\) is computed from the SST-59 kernel, and \(E_{\mathrm{bind}}^{\mathrm{bridge}}\) is the bridge functional.

\subsection{SEMF bridge functional (drop-in)}
We take
\begin{equation}
E_{\mathrm{bind}}^{\mathrm{bridge}}(Z,N)
= a_V A - a_S A^{2/3} - a_C \frac{Z(Z-1)}{A^{1/3}} - a_A\frac{(A-2Z)^2}{A} + \delta(A,Z),
\label{eq:SEMFbridge}
\end{equation}
with pairing
\[
\delta(A,Z)=
\begin{cases}
 +a_P A^{-1/2},& Z,N\ \text{even},\\
 0,& A\ \text{odd},\\
 -a_P A^{-1/2},& Z,N\ \text{odd}.
\end{cases}
\]
Equation~\eqref{eq:SEMFbridge} is standard (orthodox) nuclear phenomenology; within SST-59 it is explicitly a \emph{bridge model}, not a derivation.
Any SST-specific topological modulation (optional) must enter as a controlled deformation of the coefficients \(a_i\), leaving the structure intact.

\section{Derivation of Factors}

\subsection{Geometric Base Mass ($M_0$)}
Using the canonical SST constants for the vortex core radius $r_c$, core density $\rho_{\text{core}}$, and swirl velocity $v_{\circlearrowleft}$, the unit mass is derived as:
\begin{equation}
    M_0 = \frac{\pi r_c^3 \cdot \frac{1}{2}\rho_{\text{core}} v_{\circlearrowleft}^2}{c^2} \approx 2.277 \times 10^{-31} \text{ kg} \approx \mathbf{0.1277 \text{ MeV}/c^2}
\end{equation}
This defines the energy scale of a "naked" vortex filament of length $L=r_c$.

\subsection{Spectral Suppression ($k^{-3/2}$)}
Treating the vortex loop as a dynamical object subject to Beltrami flow conditions ($\nabla \times \mathbf{u} = \lambda \mathbf{u}$), the partition function $Z$ for quantum fluctuations is dominated by the 1-loop determinant. For a mode $k$ with eigenvalue $\lambda_k \propto k$ in 3 spatial dimensions:
\begin{equation}
    Z_k \propto (\det \mathcal{O}_k)^{-1/2} \propto (k^3)^{-1/2} = k^{-3/2}
\end{equation}
This factor suppresses the effective mass of high-frequency (complex) knot modes.

\section{Results and Validation}



\begin{table}[h]
\centering
\begin{tabular}{lccccc}
\toprule
\textbf{Particle} & \textbf{Topology} & \textbf{Geometry} ($M_{geom}$) & \textbf{Factors} & \textbf{SST Mass} & \textbf{Exp. Mass} \\
\midrule
Electron & $3_1$ (Trefoil) & 2.09 MeV & $k=3, S=0$ & \textbf{0.40 MeV} & 0.511 MeV \\
Muon & $5_1$ (Cinquefoil) & 3.38 MeV & $k=5, S \approx 0.38$ & \textbf{95.6 MeV} & 105.7 MeV \\
Proton & $3 \times (5_2/6_1)$ & $\sim 10.2$ MeV & $n=3, S=1$ & \textbf{939.8 MeV} & 938.3 MeV \\
\bottomrule
\end{tabular}
\caption{Comparison of predicted masses versus experimental values.}
\label{tab:results}
\end{table}

\subsection{The Proton-Electron Ratio}
The model successfully bridges the gap between the electron scale ($\sim 0.5$ MeV) and the nucleon scale ($\sim 1$ GeV). The proton mass is generated not by geometric length alone, but by the \textit{amplification} of the geometric mass via the vacuum coupling factor ($4/\alpha$) and the coherence factor for $n=3$ components ($3^{-1/\phi} \approx 0.5$).

\section{Conclusion}
Swirl String Theory provides a geometric basis for the mass hierarchy. The distinction between generations is identified as a topological phase transition in vacuum shielding. This framework eliminates the need for arbitrary Yukawa couplings, replacing them with discrete topological invariants ($k, n, g$) and a single fundamental mass scale $M_0$.

\begin{thebibliography}{9}
\bibitem{kelvin} Thomson, W. (1867). \textit{On Vortex Atoms}. Proc. Roy. Soc. Edin.
\bibitem{ideal_knots} Cantarella, J., et al. (2002). \textit{On the Minimum Ropelength of Knots}.
\bibitem{sst_canon} User. (2025). \textit{Swirl String Theory Canon}.

\bibitem{Lee2025DoubleTwist}
J.~Lee (2025),
\textit{Rational Concordance of Double Twist Knots},
arXiv:2504.07636v1 (Apr 11, 2025).

\bibitem{Cha2007}
J.~C.~Cha (2007),
\textit{Rational concordance and the Grope filtration of the classical knot concordance group},
Trans. Amer. Math. Soc. \textbf{360} (2007), no.~1, 405--436.

\bibitem{MY1987}
Y.~Matsumoto and S.~Yamada (1987),
\textit{Rational-slice knots},
Topology Appl. \textbf{27} (1987), no.~1, 25--35.

\bibitem{Siebenmann1975}
L.~Siebenman (1975),
\textit{Exercises sur les noeuds rationnels},
mimeographed notes.

\bibitem{Weizsaecker1935}
C.~F.~von Weizs\"acker (1935),
\textit{Zur Theorie der Kernmassen},
Z. Phys. \textbf{96}, 431--458.
\end{thebibliography}

\end{document}